\appendix


\begin{table*}[h!]\centering
    \begin{minipage}{0.98\linewidth}\vspace{0mm}    
    \centering
    \begin{tcolorbox} 
        \centering
        % \hspace{-6mm}
        \begin{tabular}{p{0.99\linewidth}}
        % \hspace{1mm}
        \begin{minipage}{0.99\linewidth}
User: The following is a story enclosed in three single quotes ```\textcolor[RGB]{243, 156, 18}{Once upon a time, there was a teddy bear named Teddy who dreamed of traveling the world. One day, his dream came true to travel around the world. Teddy sat in the airport lobby and traveled to many places of interest. Along the way, Teddy visited the Eiffel Tower, the Great Wall, and the pyramids. In Paris, Teddy had a picnic and enjoyed some delicious croissants. At the Great Wall of China, he climbed to the top and marveled at the breathtaking view. And in Egypt, he explored the pyramids and even found a secret treasure hidden inside. After his exciting journey, Teddy was eventually reunited with his owner who was thrilled to hear about all of Teddy's adventures. From that day on, Teddy always dreamed of traveling the world again and experiencing new and exciting things.}''' I would like to request your assistance in writing a script for a vlog based on this story. Provide the script in JSON format(do not answer anything else) with the following keys: video fragment id, video fragment description. You must answer like the content in the following three single quotes: ```[\{``video fragment id": 1, ``video fragment description": ``(the description, describe the characters, actions, and backgrounds in the video fragment)",\}, \{
``video fragment id": 2, ``video fragment description": ``(the description, describe the characters, actions, and backgrounds in the video fragment)",\}]'''  The descriptions of the video segments should adhere to the following guidelines: 1. Fits the original storyline. 2. All video fragment descriptions cannot conflict with each other, and the descriptions corresponding to successive fragments in the original story must have a certain continuity. 3. The description only needs to describe the visual elements presented, such as the subject, action, background, etc., and do not appear useless descriptions, such as mental activities.

GPT-4: \textcolor[RGB]{37, 204, 247}{[\{``video fragment id": 1, ``video fragment description": ``The screen opens to the sight of a well-loved teddy bear named Teddy. In the background, the hustle and bustle of an airport can be seen with people hurrying about to catch their flights."\}, \{``video fragment id": 2, ``video fragment description": ``The next scene shows Teddy looking out from the backpack, wide-eyed at the marvels of the Eiffel Tower in Paris. The iconic tower dominates the skyline, and Teddy is seen seated on a picnic blanket, with crumbs from a croissant nearby."\}, \{``video fragment id": 3, ``video fragment description": ``The setting shifts to the imposing stone structure of the Great Wall of China. Teddy is seen at the top of the wall, looking out over the breathtaking view. Views of the surrounding landscape are beautifully captured in wide shots."\}, \{``video fragment id": 4, ``video fragment description": ``In the next transition, we see Teddy exploring the mysterious pyramids of Egypt. He is captured in various poses around the ancient structures. A climactic moment occurs when Teddy discovers a secret treasure hidden deep within one of the pyramids."\}, \{``video fragment id": 5, ``video fragment description": ``The final scenes show Teddy being reunited with his owner. The joyous reunion takes place in a cozy home setting. The owner's surprised reaction upon hearing Teddy's adventures is clearly depicted. As the camera pulls away, Teddy sits happily in a familiar comforting spot, dreaming of more adventures."\}]
}
 \end{minipage}
        \end{tabular}
    \end{tcolorbox}
    \vspace{-2mm}
    \caption{\textbf{The Rough Script planning process of Vlogger.} The \textcolor[RGB]{37, 204, 247}{Rough Script} and \textcolor[RGB]{243, 156, 18}{Story} are marked by \textcolor[RGB]{37, 204, 247}{spiro disco ball} and \textcolor[RGB]{243, 156, 18}{orrange} respectively.
    }
    \label{tab:instruction_1}
\end{minipage}
\end{table*}
\begin{table*}[h!]\centering
    \begin{minipage}{0.98\linewidth}\vspace{0mm}    
    \centering
    \begin{tcolorbox} 
        \centering
        % \hspace{-6mm}
        \begin{tabular}{p{0.99\linewidth}}
        % \hspace{1mm}
        \begin{minipage}{0.99\linewidth}
User: I want to make a vlog, the following is the original story enclosed in three single quotes ```\{\textcolor[RGB]{243, 156, 18}{Story}\}''' and this is my vlog production script in the following three single quotes ```\{\textcolor{cyan}{Rough Script}\}''', a paragraph in the script corresponds to a clip of the vlog, but the description of some video clips is too complicated, please help me analyze and rewrite a video script, split each description into several short descriptions and as more as possible. For example, if there is one paragraph in the script I gave you, then you should split it into fifteen paragraphs. Provide me the answer in JSON format(do not answer anything else) with the following keys: video fragment id, video fragment description. You must answer like the content in the following three single quotes: ```[\{``video fragment id": 1, ``video fragment description": ``(the description)",\}, \{``video fragment id": 2, ``video fragment description": ``(the description)",\}, \{``video fragment id": 3, ``video fragment description": ``(the description)",\}]'''. Remember to make sure that the description of each video clip is not long, no more than ten words, but there can be so many video clips. As many fragments as possible, as detailed as possible! Ensure that all description statements are as natural and syntactically correct as possible. Very important!!!: Try to have only one character in the description and avoid complex actions in the video fragment description, such as loaded in, fight, etc.

GPT-4: \textcolor{cyan}{[\{``video fragment id": 1, ``video fragment description": ``One day, Teddy was planning his trip on paper."\}, \{``video fragment id": 2, ``video fragment description": ``Backpack is in an airport."\}, \{``video fragment id": 3, ``video fragment description": ``People rush around the airport."\}, \{``video fragment id": 4, ``video fragment description": ``Teddy peeks out from backpack."\}, \{``video fragment id": 5, ``video fragment description": ``Eiffel Tower comes into view."\}, \{``video fragment id": 6, ``video fragment description": ``Teddy sits on a blanket."\}, \{``video fragment id": 7, ``video fragment description": ``Croissant crumbs lie nearby."\}, \{``video fragment id": 8, ``video fragment description": ``Scene shifts to Great Wall."\}, \{``video fragment id": 9, ``video fragment description": ``Teddy is at the wall's top."\}, \{``video fragment id": 10, ``video fragment description": ``Wide shot of surrounding landscape."\}, \{``video fragment id": 11, ``video fragment description": ``Transition to Egyptian pyramids."\}, \{``video fragment id": 12, ``video fragment description": ``Teddy poses near the structures."\}, \{``video fragment id": 13, ``video fragment description": ``A secret treasure is discovered."\}, \{``video fragment id": 14, ``video fragment description": ``Transition to cozy home setting."\}, \{``video fragment id": 15, ``video fragment description": ``Teddy reunites with his owner."\}, \{``video fragment id": 16, ``video fragment description": ``Owner reacts with surprise."\}, \{``video fragment id": 17, ``video fragment description": ``Teddy sits in a comfortable spot."\}, \{``video fragment id": 18, ``video fragment description": ``Teddy dreams of more adventures."\}
]
}
 \end{minipage}
        \end{tabular}
    \end{tcolorbox}
    \vspace{-2mm}
    \caption{\textbf{The Detail Script planning process of Vlogger.} The \textcolor{cyan}{Detail Script}, \textcolor[RGB]{37, 204, 247}{Rough Script} and \textcolor[RGB]{243, 156, 18}{Story} are marked by \textcolor{cyan}{cyan}, \textcolor[RGB]{37, 204, 247}{spiro disco ball}, and \textcolor[RGB]{243, 156, 18}{orrange} respectively.
    }
    \label{tab:instruction_2}
\end{minipage}
\end{table*}
\begin{table*}[h!]\centering
    \begin{minipage}{0.98\linewidth}\vspace{0mm}    
    \centering
    \begin{tcolorbox} 
        \centering
        % \hspace{-6mm}
        \begin{tabular}{p{0.99\linewidth}}
        % \hspace{1mm}
        \begin{minipage}{0.99\linewidth}
User: The following is a story enclosed in three single quotes ```\{\textcolor[RGB]{243, 156, 18}{story}\}'''. I want to make a video according to this story, this is my video production script in the following three single quotes ```\{\textcolor{cyan}{Detail Script,}\}''', a paragraph in the script corresponds to a clip of the video. However, there may be some plots in the story missing, such as important plot missing, or transitions between pictures, please check and complete it for me. Provide me the answer in JSON format(do not answer anything else) with the following keys: video fragment id, video fragment description. You must answer like the content in following three single quotes: ```[\{``video fragment id": 1, ``video fragment description": ``(the description)",\}, \{ ``video fragment id": 2, ``video fragment description": ``(the description)",\}, \{``video fragment id": 3, ``video fragment description": ``(the description)",\}]'''. Remember to make sure that the description of each video clip is not long, no more than fifteen words, but there can be so many video clips. Each description should include the subject, place, and action as much as possible. As more fragment as possible, as detail as possible! Read this script carefully and don't pull down any details.

GPT-4: \textcolor[RGB]{52, 152, 219}{[\{``video fragment id": 1, ``video fragment description": ``Teddy is planning on paper."\}, \{``video fragment id": 2, ``video fragment description": ``The globe rotates with various locations in the background."\}, \{``video fragment id": 3, ``video fragment description": ``Teddy sitting in the airport lobby."\}, \{``video fragment id": 4, ``video fragment description": ``The plane took off from the runway."\}, \{``video fragment id": 5, ``video fragment description": ``Teddy arrives in Paris."\}, \{``video fragment id": 6, ``video fragment description": ``Teddy under the Eiffel Tower."\}, \{``video fragment id": 7, ``video fragment description": ``Teddy is enjoying a picnic."\}, \{``video fragment id": 8, ``video fragment description": ``Arrive at the Great Wall."\}, \{``video fragment id": 9, ``video fragment description": ``Teddy climbs the Great Wall."\}, \{``video fragment id": 10, ``video fragment description": ``Teddy watches the scenery from the Great Wall."\}, \{``video fragment id": 11, ``video fragment description": ``Teddy leaves the Great Wall."\}, \{``video fragment id": 12, ``video fragment description": ``Teddy is sitting on the bus, the scenery outside the window is beautiful."\}, \{``video fragment id": 13, ``video fragment description": ``Planes fly over the sky."\}, \{``video fragment id": 14, ``video fragment description": ``Teddy lands in a sandy landscape."\}, \{``video fragment id": 15, ``video fragment description": ``Pyramids gradually appeared in front of Teddy."\}, \{``video fragment id": 16, ``video fragment description": ``Teddy is entering the pyramid."\}, \{``video fragment id": 17, ``video fragment description": ``Teddy walks in the dark corridor."\}, \{``video fragment id": 18, ``video fragment description": ``Teddy suddenly walked into a bright room."\}, \{``video fragment id": 19, ``video fragment description": ``A treasure is revealed."\}, \{``video fragment id": 20, ``video fragment description": ``Teddy found his master with the treasure."\}, \{``video fragment id": 21, ``video fragment description": ``Room is filled with adventure mementos."\}, \{``video fragment id": 22, ``video fragment description": ``Teddy is looking at his mementos and camera rotation."\}, \{``video fragment id": 23, ``video fragment description": ``The globe gradually stopped rotating and the background gradually turned black."\}, \{``video fragment id": 24, ``video fragment description": ``Teddy sits in a chair and dreams about his next travels."\}]}

 \end{minipage}
        \end{tabular}
    \end{tcolorbox}
    \vspace{-2mm}
    \caption{\textbf{The Complete Script planning process of Vlogger.} The \textcolor[RGB]{52, 152, 219}{Complete Script}, \textcolor{cyan}{Detail Script} and \textcolor[RGB]{243, 156, 18}{Story} are marked by \textcolor[RGB]{52, 152, 219}{sea blue}, \textcolor{cyan}{cyan} and \textcolor[RGB]{243, 156, 18}{orrange} respectively.
    }
    \label{tab:instruction_3}
\end{minipage}
\end{table*}
\begin{table*}[h!]\centering
    \begin{minipage}{0.98\linewidth}\vspace{0mm}    
    \centering
    \begin{tcolorbox} 
        \centering
        % \hspace{-6mm}
        \begin{tabular}{p{0.99\linewidth}}
        % \hspace{1mm}
        \begin{minipage}{0.99\linewidth}
User: The following is a story enclosed in three single quotes ```\{\textcolor[RGB]{243, 156, 18}{story}\}''', I want to make a video according to this story, this is my video production script in the following three single quotes ```\{\textcolor[RGB]{52, 152, 219}{Complete Script}\}''', a paragraph in the script corresponds to a clip of the video, Now that you know that 16-frame videos have a length of 2 seconds, please help me plan how much time it will take for each video clip to fully interpret the meaning of the script. Provide me the answer in JSON format(do not answer anything else) with the following keys: video fragment id, time. You must answer like the content in the following three single quotes: [\{``video fragment id": 1, ``time": 2,\}, \{``video fragment id": 2, ``time": 10,\}, \{``video fragment id": 3, ``time": 6,\}]

GPT-4: \textcolor[RGB]{30, 144, 255}{[\{``video fragment id": 1, ``time": 6\}, \{``video fragment id": 2, ``time": 10\}, \{``video fragment id": 3, ``time": 6\}, \{``video fragment id": 4, ``time": 7\}, \{``video fragment id": 5, ``time": 8\}, \{``video fragment id": 6, ``time": 6\}, \{``video fragment id": 7, ``time": 8\}, \{``video fragment id": 8, ``time": 9\}, \{``video fragment id": 9, ``time": 7\}, \{``video fragment id": 10, ``time": 8\}, \{``video fragment id": 11, ``time": 7\}, \{``video fragment id": 12, ``time": 10\}, \{``video fragment id": 13, ``time": 7\}, \{``video fragment id": 14, ``time": 6\}, \{``video fragment id": 15, ``time": 8\}, \{``video fragment id": 16, ``time": 5\}, \{``video fragment id": 17, ``time": 8\}, \{``video fragment id": 18, ``time": 6\}, \{``video fragment id": 19, ``time": 9\}, \{``video fragment id": 20, ``time": 6\}, \{``video fragment id": 21, ``time": 7\}, \{``video fragment id": 22, ``time": 8\}, \{``video fragment id": 23, ``time": 10\}, \{``video fragment id": 24, ``time": 6\}]}

 \end{minipage}
        \end{tabular}
    \end{tcolorbox}
    \vspace{-2mm}
    \caption{\textbf{The Script planning process of Vlogger.} The \textcolor[RGB]{30, 144, 255}{Script}, \textcolor[RGB]{52, 152, 219}{Complete Script} and \textcolor[RGB]{243, 156, 18}{Story} are marked by \textcolor[RGB]{30, 144, 255}{clear shill}, \textcolor[RGB]{52, 152, 219}{sea blue} and \textcolor[RGB]{243, 156, 18}{orrange} respectively.
    }
    \label{tab:instruction_4}
\end{minipage}
\end{table*}
\begin{table*}[t]
    % \scriptsize
    \centering
    \setlength\tabcolsep{4.5pt}
    \resizebox{1\linewidth}{!}{%
    \begin{tabular}{l | c | c | c | c | c | c | c | c | c | c | c | c | c | c | c}
        \toprule
    \multicolumn{1}{c |}{ } & \multicolumn{3}{| c}{\textbf{Teddy Travel}} & \multicolumn{3}{| c}{\textbf{Alfie Work}} & \multicolumn{3}{| c}{\textbf{Tiger Moon}} & \multicolumn{3}{| c}{\textbf{Pete Skiing}} & \multicolumn{3}{| c}{\textbf{Anna Cooking}} \\
    \midrule
        \textbf{Method} & \textbf{CI ($\uparrow$)} & \textbf{CT ($\uparrow$)} & \textbf{VD ($\uparrow$)} & \textbf{CI ($\uparrow$)} & \textbf{CT ($\uparrow$)} & \textbf{VD ($\uparrow$)} & \textbf{CI ($\uparrow$)} & \textbf{CT ($\uparrow$)} & \textbf{VD ($\uparrow$)} & \textbf{CI ($\uparrow$)} & \textbf{CT ($\uparrow$)} & \textbf{VD ($\uparrow$)} & \textbf{CI ($\uparrow$)} & \textbf{CT ($\uparrow$)} & \textbf{VD ($\uparrow$)}  \\
        
        \midrule
        w/o \small{autoregressive}~\cite{Zhu2023MovieFactoryAM} & 0.5659 & 0.2846 & 48s & 0.5537 & 0.2641 & 1min20s & 0.6918 & 0.2819 & 1min10s & 0.4946 & 0.2627 & 54s & 0.5356 & 0.288 & 1min04s \\
        
        only \small{autoregressive}~\cite{villegas2023phenaki} & 0.5578 & 0.2554 & 2min47s & 0.5445 & 0.2361 & 5min12s & 0.6859 & 0.262 & 4min17s & 0.5033 & 0.2463 & 3min22s & 0.5295 & 0.2675 & 4min11s \\
        
        \textbf{Ours} & \textbf{0.653} & \textbf{0.2847} & \textbf{2min47s} & \textbf{0.6168} & \textbf{0.262} & \textbf{5min12s} & \textbf{0.7848} & \textbf{0.2802} & \textbf{4min17s} & \textbf{0.5393} & \textbf{0.2639} & \textbf{3min22s} & \textbf{0.553} & \textbf{0.2871} & \textbf{4min11s} \\
        \bottomrule
        % \vspace{-1.75em}
    \end{tabular}
    }
    \caption{\textbf{Detail ablation for generation process}.
    CI, CT, and VD refer to CLIP-I, CLIP-T, and Vlog Duration respectively.
    }
    \vspace{-1em}
    \label{tab:framework_ablation_detail}
\end{table*}
\begin{table*}[t]
\begin{minipage}[t]{1\linewidth}
    \centering
    \setlength\tabcolsep{6pt}
    \resizebox{1\linewidth}{!}{%
    \begin{tabular}{c | c | c | c | c | c | c | c | c | c | c | c | c | c | c | c | c | c | c | c | c | c }
    \toprule
    \multicolumn{2}{c |}{ } & \multicolumn{4}{| c}{\textbf{Experimental 1003}} & \multicolumn{4}{| c}{\textbf{Music 1179}} & \multicolumn{4}{| c}{\textbf{Animation 1175}} & \multicolumn{4}{| c}{\textbf{Sport 905}} & \multicolumn{4}{| c}{\textbf{Ads and Commercials 1229}} \\
    \midrule
        \textbf{TTCA} & \textbf{MTP} & 0/16 () & 1/15 () & 3/13 () & 5/11 () & 0/16 () & 1/15 () & 3/13 () & 5/11 () & 0/16 () & 1/15 () & 3/13 () & 5/11 () & 0/16 () & 1/15 () & 3/13 () & 5/11 () & 0/16 () & 1/15 () & 3/13 () & 5/11 ()  \\
        \midrule
        \xmark & \xmark & 326.75 & 244.45 & 302.43 & 260.52 & 362.06 & 254.73 & 377.09 & 318.92 & 283.49 & 204.61 & 272.61 & 277.08 & 608.13 & 332.09 & 433.65 & 313.92 & 339.66 & 190.68 & 268.9 & 211.51 \\
        
        \cmark & \xmark & 319.93 & 197.03 & 251.3 & 226.34 & 357.28 & 206.5 & 261.78 & 221.05 & 252.42 & 174.8 & 262.87 & 243.34 & 556.09 & 261.2 & 315.4 & 218.27 & 332.7 & 177.28 & 205.07 & 166.19 \\
        
        \cmark & \cmark & \textbf{227.3} & \textbf{124.23} & \textbf{113.13} & \textbf{122.28} & \textbf{250.4} & \textbf{132.16} & \textbf{130.2} & \textbf{128.78} & \textbf{207.3} & \textbf{115.3} & \textbf{130.14} & \textbf{104.94} & \textbf{427.79} & \textbf{196.11} & \textbf{181.84} & \textbf{148.87} & \textbf{239.88} & \textbf{107.05} & \textbf{101.6} & \textbf{109.17} \\
        \bottomrule
        % \vspace{-1.75em}
    \end{tabular}
    }
    \subcaption{Ablation on Experimental, Music, Animation, Sport, and Ads and Commercials categories of Vimeo11K.}
    \end{minipage}

\begin{minipage}[t]{1\linewidth}
    \centering
    \setlength\tabcolsep{6pt}
    \resizebox{1\linewidth}{!}{%
    \begin{tabular}{c | c | c | c | c | c | c | c | c | c | c | c | c | c | c | c | c | c | c | c | c | c }
    \toprule
    \multicolumn{2}{c |}{ } & \multicolumn{4}{| c}{\textbf{Travel 1288}} & \multicolumn{4}{| c}{\textbf{Branded Content 1173}} & \multicolumn{4}{| c}{\textbf{Narrative 1010}} & \multicolumn{4}{| c}{\textbf{Comedy 1267}} & \multicolumn{4}{| c}{\textbf{Documentary 1064}} \\
    \midrule
        \textbf{TTCA} & \textbf{MTP} & 0/16 () & 1/15 () & 3/13 () & 5/11 () & 0/16 () & 1/15 () & 3/13 () & 5/11 () & 0/16 () & 1/15 () & 3/13 () & 5/11 () & 0/16 () & 1/15 () & 3/13 () & 5/11 () & 0/16 () & 1/15 () & 3/13 () & 5/11 ()  \\
        \midrule
        \xmark & \xmark & 432.36 & 215.06 & 285.86 & 244.87 & 246.83 & 174.44 & 221.33 & 188.21 & 301.44 & 202.36 & 310.94 & 232.17 & 244.46 & 172.31 & 260.8 & 219.12 & 338.39 & 183.15 & 240.09 & 199.32 \\
        
        \cmark & \xmark & 382.39 & 178.86 & 220.89 & 172.71 & 249.88 & 135.09 & 173.19 & 150.78 & 272.07 & 173.55 & 224.69 & 192.73 & 232.65 & 135.15 & 191.19 & 182.11 & 380.86 & 143.28 & 199.26 & 161.65 \\
        
        \cmark & \cmark & \textbf{283.64} & \textbf{123.55} & \textbf{116.24} & \textbf{92.99} & \textbf{209.6} & \textbf{98.77} & \textbf{98.19} & \textbf{88.17} & \textbf{242.63} & \textbf{135.59} & \textbf{107.18} & \textbf{109.89} & \textbf{205.04} & \textbf{94.03} & \textbf{103.85} & \textbf{95.3} & \textbf{283.38} & \textbf{108.28} & \textbf{97.88} & \textbf{92.69} \\
        \bottomrule
        % \vspace{-1.75em}
    \end{tabular}
    }
    \subcaption{Ablation on Travel, Branded Content, Narrative, Comedy, and Documentary categories of Vimeo11K.}
    \end{minipage}
    
    \caption{\textbf{Detail ablation for temporal text cross attention and mixed training paradigm. }. Detailed experimental data on 10 major categories of Vimeo11K.
    }
    \vspace{-0.5cm}
    \label{tab:vimeo_detail}
\end{table*}

% \setcounter{section}{0}
\section{Top-Down Planning}
\label{sub:planning}

As depicted in Tab. \ref{tab:instruction_1}, \ref{tab:instruction_2}, \ref{tab:instruction_3} and \ref{tab:instruction_4}, we encourage users to follow the provided example as guidance. 
These prompts instruct the LLM to generate structured outputs, specifically in the form of JSON code snippets. 
Concurrently, the examples furnished in these tables qualitatively emphasize the essentiality of a progressive approach in script generation, thereby rendering them qualitatively insightful.

\section{ShowMaker Implementation Details}
\textbf{Stage-1 Settings}.
% For architecture, we expand the input channel of the conv-in layer of our base video pre-train model \cite{Wang2023LAVIEHV} from 4 to 9. 
% Additionally, after we incorporate temporal text cross attention into the U-Net architecture with a total parameter count of 1.3B. 
% For training, we train the entire U-Net, only excluding the spatial image cross attention parts.
In terms of architecture, we modify the input channel of the convolutional input layer in our base video pre-training model \cite{Wang2023LAVIEHV}, expanding it from 4 to 9. 
Furthermore, the integration of temporal text cross attention into the U-Net architecture results in a total parameter count of 1.3 billion. 
In terms of training, we train the entire U-Net, excluding only the spatial image cross attention components.
We set the value of the mixed training paradigm $\alpha$ and $m$ to 0.6 and 6 respectively. 
We establish the probability of dropping both the textual prompt and visual prompt at 0.1.
The learning rate we employ is 1e-4, and the batch size is 384, comprising a 16-frame video and 6 images per batch. 
We perform 54,000 steps total, requiring approximately five days to complete. 
% With the training conducted on 128 A100 units, we perform 54000 steps, requiring approximately five days to complete. 
It is worth noting that when training image data, the mask probability is set to 1, and the image features do not undergo temporal self attention and temporal text cross attention module.

\noindent\textbf{Stage-2 Settings}.
For architecture, we incorporate OpenCLIP ViT-H/14 \cite{ilharco2021openclip} as our image encoder. We set the spatial image cross attention $\beta$ to 1 during training.
For training, the U-Net is frozen while only permitting the training of spatial image cross attention components, thereby limiting the number of trainable parameters to only 22M. 
We follow a joint approach involving videos and images. 
However, unlike the first stage, we only need to train the video generation pattern in the second stage. 
During the training of spatial image cross attention, besides using a textual prompt as a condition, a visual prompt is also needed. 
The reference image for a given image is the image itself, while for videos, a random frame is selected from it. 
The probability of dropping both the textual prompt and visual prompt is set to 0.1. 
The learning rate used is 1e-4, and the batch size is 256, comprising a 16-frame video and 6 images per batch. 
We perform 62,000 steps, requiring approximately five days to complete. 

For both two stages, the AdamW \cite{Loshchilov2017DecoupledWD} optimizer is consistently employed. 
All images and videos used during training are subjected to online resizing, yielding a resolution of $320$$\times$$520$ resolution.
Regarding the video sampling process within the WebVid10M \cite{bain2021frozen}, we adopt a random sampling technique with an interval of 6 frames.





\section{Comparison with State-of-the-art}

\textbf{UCF-101} \cite{Soomro2012UCF101AD}.
To ensure accurate and stable computation of the FVD, we use the same sampling strategy as \cite{Wang2023LAVIEHV}, where a hundred videos are sampled for each class or hand-crafted prompt.
We first generate the video at a resolution of $320$$\times$$512$ and then resize it to match the $240$$\times$$320$ resolution of UCF-101.
Initially, videos are generated at a resolution of $320$$\times$$512$. 
Subsequently, these videos are resized to align with the $240$$\times$$320$ resolution of the UCF-101 dataset.
For the calculation of FVD \cite{unterthiner2019accurate}, we utilize a codebase derived from \cite{skorokhodov2022styleganv}.
\noindent\textbf{Kinetics-400} \cite{Kay2017TheKH}.
For Text-to-Video (T2V) generation within the Kinetics-400 dataset, we employ GPT-4 \cite{OpenAI2023GPT4TR} to construct a hand-crafted prompt that corresponds with each class label. 
Given the variable video resolutions present in the Kinetics-400 dataset, the ground truth videos are resized to a $320$$\times$$512$ resolution to match our generated videos.
\textbf{MSR-VTT} \cite{Chen2021TheMT}.
The MSR-VTT test set comprises 2,990 examples, each accompanied by 20 descriptions. 
For each example, we generate a video (comprising 16 frames at $320$$\times$$512$ resolution) using one randomly selected prompt. 
This procedure results in the creation of 2,990 videos.
The CLIPSIM \cite{Wu2021GODIVAGO} and FID \cite{parmar2022aliased} for the 47,840 frames are then computed following the methodology in \cite{singer2023makeavideo}.
\textbf{UCF-101 1000 Frames}.
In alignment with TATS \cite{Ge2022LongVG}, we use the hand-crafted prompts of UCF-101 as \cite{Ge2023PreserveYO} to generate videos and partition the generated long video into clips, each consisting of 16 frames. 
The FVD is computed between these clips and the ground truth.


\section{Ablation Study}
\label{sub:ablation}

\subsection{Vlog Generation Process}
As detailed in Tab. \ref{tab:framework_ablation_detail}, we used GPT-4 \cite{OpenAI2023GPT4TR} to automatically generate five scripts, namely \textit{Teddy Travel}, \textit{Alfie Work}, \textit{Tiger Moon}, \textit{Pete Skiing} and \textit{Anna Cooking} respectively.
Then we uniformly use the planning process of Vlogger to get the Script.
However, although we use the same Script, MovieFactory \cite{Zhu2023MovieFactoryAM} and Phenaki \cite{villegas2023phenaki} cannot use visual prompts due to their framework design, and can only use textual prompts for video generation.


\subsection{Temporal Text Cross Attention (TTCA) and Mixed Training Paradigm (MTP)}
As detailed in Tab. \ref{tab:vimeo_detail}, we get the video data of Vimeo11K from \href{https://vimeo.com/}{vimeo.com} and then use VideoChat \cite{Li2023VideoChatCV} to caption the videos.
Vimeo11K mainly contains 10 major categories of videos, namely \textit{Experimental}, \textit{Music}, \textit{Animation}, \textit{Sports}, \textit{Ads and Commercials}, \textit{Travel}, \textit{Branded Content}, \textit{Narrative}, \textit{Comedy}, and \textit{Documentary} respectively.
The number of videos included in each of the 10 categories is 1003, 1079, 1175, 905, 1229, 1288, 1173, 1010, 1267, and 1064 respectively, for a total of 11293.
The dataset and code for the test content will be released afterward. 


\section{Visualization}
\label{sub:visualization}
We place the visualization results generated by our Vlogger in the supplementary material. 
Under the premise of not disrupting the files within the folder, please click on the following link: \href{https://zhuangshaobin.github.io/Vlogger.github.io/}{Project Page}.

